\section{What is proved, checked, and assumed}
\label{sec:appendix:limits}

\tab\ref{tab:claims} in \S\ref{sec:formal-status} gives the short version.  This appendix gives the
long one, claim by claim, because the differences between "proved on paper", "machine-checked in
Lean" and "held by a gate" are exactly what a reader of a systems paper should be able to audit.

\subsection{Safety of the region machine}
\label{sec:appendix:safe}

\thm\ref{thm:safe} is proved in \S\ref{sec:thm-safe} from four invariants: the root table of a region
covers every live slot reachable from outside it (I1), every live slot that holds a reference points
at a live slot (I2), no live slot points at released storage (I3), and the shade set over-approximates
the reachable set at the moment marking began (I4).  The four release layers are the four ways a slot
can become free, and the proof is a case split over them, which is why it is worth naming the Lean
counterpart of each case:
\begin{itemize}
\item \emph{block exit} — \code{ExtCBlock.safe\_freeBlk}, with \code{referrer\_survives\_free}
covering the case in which a reference outlives the block that allocated the storage;
\item \emph{static placement} — \code{ExtCBlock.safe\_promote}, which in this design is the
statement that a value allocated directly in a shallower region is never the target of a pointer to
the region it escaped from;
\item \emph{sweep} — \code{ExtCRegion.World.safe\_of\_invariants} (given an over-approximation of
the live set) and \code{ExtCReclaim.safe\_kill\_unreachable} (a region that is unreachable from
every seed may be freed as a unit);
\item \emph{release at the orphaning store} — \code{ExtCBlock.referrer\_survives\_free} together
with the four static side conditions of the rule, which make the freed set a subset of what a full
sweep would free.
\end{itemize}
The module \code{ExtCRegion} is the largest in the development (742 of the 2,555 lines) because it
also carries the negative results of \S\ref{sec:appendix:necessity}.

\subsection{Why the assumptions are not decorative}
\label{sec:appendix:necessity}

An invariant that is assumed in a safety theorem should be shown to be necessary, or the theorem
says less than it appears to.  The Lean development therefore proves, for three of the four
invariants, that dropping it makes the safety statement false:
\begin{enumerate}
\item \code{ExtCRegion.ForeignRegion.necessity\_CC} exhibits two regions, each of which is locally
consistent, whose union has a dangling reference; the witness is exactly a cross-region reference
that the foreign table does not record.  This is why the compiler emits a cross-region bit for every
stored edge whose target lies in another region, and why the bit is one bit and not a count: the
proof only needs to know that the region is an entry point.
\item \code{ExtCRegion.Undershoot.necessity\_INVH} exhibits a machine in which the shade set omits a
slot that was live when marking began, and a marking pass that then frees it.  This is the content of
invariant I4, and it is why the barrier shades both the new value and the old value (SATB): shading
only one of them undershoots.
\item \code{ExtCRegion.LooseKill.necessity\_kill\_closed} exhibits a region whose kill is justified
by reachability alone while a live slot in another region still points into it; the proof needs the
closure condition that no such slot exists, which the runtime discharges by treating any registered
cross-region entry as a seed.
\end{enumerate}
Two further results delimit what can be achieved rather than justifying an assumption:
\code{ExtCPrecision.any\_sound\_rule\_le\_maxSound} shows that the store rule we emit is the most
permissive sound rule expressible in the analysis' abstraction, and
\code{ExtCModel.no\_sound\_finite\_fallback} shows that no finite fallback table can be sound for the
unknown-host case, which is why the implementation pins such hosts instead of guessing.

\subsection{Where the Lean development stops}
\label{sec:appendix:gaps}

Four honest gaps, in decreasing order of what a reader might expect them to cover:
\begin{enumerate}
\item \emph{The memory theorem is covered at the level of its law, not its statement.}
\code{ExtCGap} and \code{ExtCMemGap} establish the gap law (with a bounded allocation rate the heap
stays within a constant factor of the live set) and the monotonicity facts that
\thm\ref{thm:mem} uses, not the theorem about this runtime's trigger, its scale factor and its
wall.  Those are checked by the boundedness measurement of \tab\ref{tab:bounded} and by the
whole-program instruction budget, not by a proof.
\item \emph{The pause theorem has no Lean counterpart.}  \thm\ref{thm:pause} is a statement about
emitted code — how many words a step scans, how often it reads the clock, how much sweep credit it
can spend — and the Lean development models the region machine, not the C.  What exists is
\code{ExtCPrecision}'s footprint results for a single slice, which is one term of the bound.  The
rest is held by \tab\ref{tab:budgets} as compile-time constants and by a gate that fails when a step
walks more than its cap, plus the measured distribution of \tab\ref{tab:pause}.
\item \emph{The amortised theorem has no Lean counterpart.}  \thm\ref{thm:amortized} is a counting
argument over transitions; its implementation check is the instruction-count budget, which fails the
build when a shape's emitted code exceeds a frozen number.
\item \emph{The escape analysis is assumed, not proved.}  Assumption~\ref{as:sound} (i) is a
dataflow analysis with several hundred lines of checker code behind it.  We test it on a corpus of
148 escaped-reference programs and on every test in the tree, and the placement rule it feeds is
proved sound if the analysis is; what is not proved is the analysis itself.
\end{enumerate}
Two smaller boundaries are worth recording.  The Lean modules model reclamation as a synchronous
system, so the concurrency story of \S\ref{sec:limitations} is outside them; and the development is
checked by a script (\code{proofs/verify.sh}) that compiles the modules, scans the source for
\code{sorry} and for user-declared axioms, checks the axiom dependencies of the headline theorems,
and cross-checks a family of small cases with Z3 — so "no \code{sorry}" here means the script's
check, and the script is in the artifact.  The script compiles the modules in a fixed order and so
assumes their \code{.olean} files exist; from a clean tree they must be built in dependency order
first, after which all thirteen compile and the audit reports only \code{propext} and
\code{Quot.sound} as axioms.  The Z3 step additionally needs the \code{z3-solver} Python package;
where it is absent the remaining checks run and that one is reported as failed rather than skipped.

\subsection{What the numbers are, epistemically}

Nothing in \S\ref{sec:eval} is a theorem, and the paper does not use the words "verified" or "proved"
about a measurement.  \tab\ref{tab:crosslang}, \tab\ref{tab:pause} and \tab\ref{tab:bounded} are the
output of scripts that are in the artifact, on one machine, with the protocol of
\S\ref{sec:eval-setup}.  The properties that separate this work from a collection of measurements are
the theorems; the measurements are how we know the theorems are about a real program, and the gates
of \S\ref{sec:appendix:artifact} are how we know the program has not drifted away from them.
